% Draft: rewritten 2026-11-23 for the stocking-standard surface.
\section{Conclusion}
\label{sec:conclusion}

We analysed forest policy under fire risk as a two-layer problem on a real timber supply area: a
policy layer sets the harvest level and the stocking standard, and an enterprise responds with a
harvest schedule. The estate is the compiled TSA~24 inventory; fire is the recorded 1950--2024
regime, and the standard changes the severity of a fire that reaches a stand.

The compiled estate sustains about \SI{51.9}{\mega\cubic\meter} per period, a mean of
\SI{5.19}{\mega\cubic\meter\per\year}, and that schedule is worth \SI{8779}{million\CAD} of net
present value at a \SI{3}{\percent} real discount rate. The level sits about \SI{25}{\percent}
below the official cut; the gap is consistent with a flat harvest profile on a non-regulated age
structure, not with a missing land base.

The stocking standard is a policy lever. The 26 percent standard adds
\SIrange{32}{53}{\mega\cubic\meter} with no loss in net present value; the 55 and 84 percent
standards cost about \SI{1.2}{\percent} and \SI{1.8}{\percent} of net present value. The
composition-aware fuel projection is consistent with the operational VRI layer, and the severity
gradient is non-linear: under the calibrated mixedwood threshold only the 84 percent standard
becomes the deciduous fuel type.

The risk result is a priced hedge, not a free improvement. The 84 percent standard cuts CVaR20 by
\SI{8.1}{\mega\CAD} at a production cost of \SI{153.8}{\mega\CAD}, an exchange rate of about
\SI{19}{\mega\CAD} per \SI{1}{\mega\CAD} of tail reduction. The stand-level severity multiplier
is 0.267, but the estate CVaR20 falls only \SI{2.7}{\percent}, because the unmanaged half of the
estate keeps its fuel, conversion is gradual, and the NPV loss is dominated by early, discounted
fires. A risk-neutral planner chooses the 26 percent standard; a planner whose marginal value of
tail risk exceeds about 19 to 1 chooses the 84 percent standard.

The open question is spatial: where to place broadleaf stands as fire breaks, and how far that can
go before production suffers. Answering it requires a placement rule and a spread model; this paper
provides the production and severity surface that such an analysis would build on.
